
\documentclass[aspectratio=169,11pt]{beamer}

% Compile with XeLaTeX
\usepackage{fontspec}
\setsansfont{DejaVu Sans}
\setmonofont{DejaVu Sans Mono}
\usepackage{polyglossia}
\setdefaultlanguage{vietnamese}
\usepackage{booktabs}
\usepackage{array}
\usepackage{tabularx}
\usepackage{adjustbox}
\usepackage{listings}
\usepackage{xcolor}
\usepackage{ragged2e}
\usepackage{graphicx}
\usepackage{hyperref}

\graphicspath{{images/}{./}{/mnt/data/images/}}
\hypersetup{unicode=true}

\usetheme{Madrid}
\usecolortheme{seahorse}
\setbeamertemplate{navigation symbols}{}
\setbeamertemplate{footline}[frame number]
\setbeamertemplate{itemize items}[circle]
\setbeamertemplate{enumerate items}[default]

\definecolor{codebg}{RGB}{248,248,248}
\definecolor{codekw}{RGB}{0,72,160}
\definecolor{codestr}{RGB}{150,60,30}
\definecolor{accent}{RGB}{0,90,150}
\definecolor{placeholderbg}{RGB}{245,248,252}
\definecolor{placeholderline}{RGB}{90,120,150}

\lstset{
  basicstyle=\ttfamily\scriptsize,
  keywordstyle=\color{codekw}\bfseries,
  stringstyle=\color{codestr},
  backgroundcolor=\color{codebg},
  frame=single,
  breaklines=true,
  columns=fullflexible,
  showstringspaces=false,
  keepspaces=true,
  language=SQL
}

\newcommand{\keyword}[1]{\textcolor{accent}{\textbf{#1}}}
\newcommand{\kw}[1]{\keyword{#1}}
\newcommand{\tightlist}{\setlength{\itemsep}{2pt}\setlength{\parskip}{0pt}}
\newcommand{\smallitem}{\setlength{\itemsep}{2pt}\setlength{\parskip}{0pt}}
\newcommand{\qoption}[2]{\item[#1.] #2}

% Image placeholder: if the image exists, it is inserted; otherwise a box is shown.
% To insert real images later, place PNG files into the images/ folder with the exact filenames.
\newcommand{\imageplaceholder}[2]{%
\IfFileExists{#1}{\includegraphics[width=0.78\textwidth,height=0.58\textheight,keepaspectratio]{#1}}{%
\IfFileExists{images/#1}{\includegraphics[width=0.78\textwidth,height=0.58\textheight,keepaspectratio]{images/#1}}{%
\fcolorbox{placeholderline}{placeholderbg}{\parbox[c][0.42\textheight][c]{0.78\textwidth}{\centering\Large Chưa tìm thấy ảnh minh họa\\[2mm]\normalsize #2\\[2mm]\scriptsize Kiểm tra file: \texttt{images/\detokenize{#1}}}}}}}

\title[ER Model trong DBMS]{Giới thiệu ER Model trong DBMS}
\subtitle{Entity, Attribute, Relationship, Degree, Cardinality và Participation Constraint}
\author{Tài liệu bài giảng}
\institute{Môn học: Cơ sở dữ liệu}
\date{Cập nhật: 02/04/2026}

\AtBeginSection[]{
  \begin{frame}
    \centering
    \vfill
    {\Huge \insertsection}\par
    \vfill
  \end{frame}
}

\begin{document}

%------------------------------------------------
\begin{frame}
  \titlepage
\end{frame}

%------------------------------------------------
\begin{frame}{Mục tiêu bài giảng}
\small
Sau khi hoàn thành bài học này, người học có thể:
\begin{enumerate}
\smallitem
  \item Giải thích khái niệm \kw{Entity-Relationship Model} trong thiết kế cơ sở dữ liệu.
  \item Trình bày vai trò của \kw{ER Model} và \kw{ER Diagram}.
  \item Phân biệt các thành phần chính: \kw{entity}, \kw{attribute} và \kw{relationship}.
  \item Nhận biết \kw{strong entity} và \kw{weak entity}.
  \item Phân biệt các loại attribute: simple, key, composite, multivalued và derived.
  \item Giải thích \kw{degree}, \kw{cardinality} và \kw{participation constraint}.
  \item Vận dụng quy trình cơ bản để vẽ ER Diagram cho một bài toán thực tế.
\end{enumerate}
\end{frame}

%------------------------------------------------
\begin{frame}{Cấu trúc bài giảng}
\small
\begin{columns}[T]
\begin{column}{0.48\textwidth}
\begin{enumerate}
\smallitem
  \item Tổng quan ER Model
  \item Vai trò trong thiết kế CSDL
  \item Thành phần ER Model
  \item Entity, strong entity và weak entity
  \item Attribute và các loại attribute
\end{enumerate}
\end{column}
\begin{column}{0.48\textwidth}
\begin{enumerate}
\setcounter{enumi}{5}
\smallitem
  \item Relationship và degree
  \item Cardinality
  \item Participation constraint
  \item Cách vẽ ER Diagram
  \item Ví dụ, câu hỏi và bài tập
\end{enumerate}
\end{column}
\end{columns}
\end{frame}

%================================================
\section{Tổng quan ER Model}

%------------------------------------------------
\begin{frame}{ER Model là gì?}
\small
\kw{Entity-Relationship Model}, hay \kw{ER Model}, là mô hình khái niệm dùng để thiết kế cơ sở dữ liệu.

\vspace{2mm}
Mô hình này biểu diễn cấu trúc logic của cơ sở dữ liệu thông qua:
\begin{itemize}
\smallitem
  \item Các \kw{thực thể} cần lưu trữ dữ liệu.
  \item Các \kw{thuộc tính} mô tả thực thể.
  \item Các \kw{mối quan hệ} giữa các thực thể.
\end{itemize}
\end{frame}

%------------------------------------------------
\begin{frame}{ER Model giúp trả lời những câu hỏi nào?}
\small
ER Model giúp người thiết kế trả lời:
\begin{itemize}
\smallitem
  \item Hệ thống cần lưu trữ những đối tượng nào?
  \item Mỗi đối tượng có những thông tin gì?
  \item Các đối tượng liên hệ với nhau như thế nào?
  \item Những ràng buộc nào cần được thể hiện trong mô hình dữ liệu?
\end{itemize}

\begin{block}{ER Diagram}
Biểu diễn đồ họa của ER Model được gọi là \kw{Entity-Relationship Diagram}, viết tắt là \kw{ERD}.
\end{block}
\end{frame}

%------------------------------------------------
\begin{frame}{Hình minh họa: Tổng quan ER Model}
\centering
\imageplaceholder{introduction-to-er-model.png}{Tổng quan Entity-Relationship Model}

\vspace{2mm}
{\small Hình 1. Tổng quan về ER Model.}
\end{frame}

%------------------------------------------------
\begin{frame}{Vai trò của ER Model trong thiết kế CSDL}
\small
Quá trình thiết kế cơ sở dữ liệu thường gồm ba giai đoạn:
\begin{enumerate}
\smallitem
  \item \kw{Thu thập yêu cầu}: hỏi người dùng, khách hàng hoặc stakeholder về chức năng và dữ liệu cần quản lý.
  \item \kw{Thiết kế khái niệm hoặc logic}: ER Model mô tả entity, attribute và relationship.
  \item \kw{Thiết kế vật lý}: xác định bảng, khóa, chỉ mục, view và chi tiết triển khai trên DBMS cụ thể.
\end{enumerate}
\end{frame}

%------------------------------------------------
\begin{frame}{Vì sao ER Diagram hữu ích?}
\small
ER Diagram hữu ích vì:
\begin{itemize}
\smallitem
  \item Biểu diễn dữ liệu và quan hệ dữ liệu bằng hình ảnh.
  \item Dễ chuyển đổi sang các bảng quan hệ.
  \item Mô hình hóa được đối tượng trong thế giới thực.
  \item Không yêu cầu người xem phải hiểu sâu về DBMS cụ thể.
  \item Giúp hệ thống phức tạp trở nên dễ hiểu hơn.
\end{itemize}
\end{frame}

%------------------------------------------------
\begin{frame}{Quiz nhanh 1: Tổng quan ER Model}
\small
\textbf{Câu 1.} ER Model chủ yếu được dùng để làm gì?
\begin{enumerate}
\tightlist
  \qoption{A}{Thiết kế mô hình khái niệm của cơ sở dữ liệu}
  \qoption{B}{Tăng tốc CPU của máy chủ}
  \qoption{C}{Thiết kế giao diện người dùng}
  \qoption{D}{Nén dữ liệu ảnh}
\end{enumerate}

\textbf{Câu 2.} ER Diagram là gì?
\begin{enumerate}
\tightlist
  \qoption{A}{Một công cụ sao lưu dữ liệu}
  \qoption{B}{Biểu diễn đồ họa của ER Model}
  \qoption{C}{Một hệ điều hành cơ sở dữ liệu}
  \qoption{D}{Một ngôn ngữ lập trình web}
\end{enumerate}
\end{frame}

%------------------------------------------------
\begin{frame}{Quiz nhanh 1: Tổng quan ER Model (tiếp)}
\small
\textbf{Câu 3.} ER Model thường được dùng nhiều nhất ở bước nào trong thiết kế cơ sở dữ liệu?
\begin{enumerate}
\tightlist
  \qoption{A}{Cấu hình card mạng}
  \qoption{B}{Biên dịch mã nguồn}
  \qoption{C}{Thiết kế khái niệm hoặc logic}
  \qoption{D}{Xóa dữ liệu cũ}
\end{enumerate}
\end{frame}

%================================================
\section{Các thành phần chính của ER Model}

%------------------------------------------------
\begin{frame}{Ba thành phần cốt lõi}
\small
ER Model gồm ba thành phần cốt lõi:
\vspace{2mm}

\begin{tabularx}{\textwidth}{>{\bfseries}p{0.22\textwidth} X p{0.28\textwidth}}
\toprule
Thành phần & Ý nghĩa & Ký hiệu thường dùng trong ERD \\
\midrule
Entity & Đối tượng hoặc khái niệm cần lưu trữ dữ liệu & Hình chữ nhật \\
Attribute & Thuộc tính mô tả entity & Hình oval \\
Relationship & Mối liên kết giữa các entity & Hình thoi \\
\bottomrule
\end{tabularx}
\end{frame}

%------------------------------------------------
\begin{frame}{Hình minh họa: Components of ER Diagram}
\centering
\imageplaceholder{components-of-er-diagram.png}{Các thành phần cơ bản trong ER Diagram}

\vspace{2mm}
{\small Hình 2. Entity, Attribute và Relationship trong ER Diagram.}
\end{frame}

%================================================
\section{Entity trong ER Model}

%------------------------------------------------
\begin{frame}{Entity là gì?}
\small
\kw{Entity} là một đối tượng, sự vật, khái niệm hoặc sự kiện trong thế giới thực mà hệ thống cần lưu trữ thông tin.

\vspace{2mm}
Ví dụ:
\begin{columns}[T]
\begin{column}{0.48\textwidth}
\begin{itemize}
\smallitem
  \item \texttt{Student}
  \item \texttt{Course}
  \item \texttt{Employee}
  \item \texttt{Company}
\end{itemize}
\end{column}
\begin{column}{0.48\textwidth}
\begin{itemize}
\smallitem
  \item \texttt{Product}
  \item \texttt{Reservation}
  \item \texttt{Device}
  \item \texttt{Document}
\end{itemize}
\end{column}
\end{columns}

\vspace{2mm}
Trong cơ sở dữ liệu quan hệ, entity thường được chuyển thành \kw{bảng}.
\end{frame}

%------------------------------------------------
\begin{frame}{Entity Type, Entity Instance và Entity Set}
\small
\begin{block}{Entity Type}
Mô tả cấu trúc chung của một loại entity, ví dụ: \texttt{Student(student\_id, full\_name, date\_of\_birth, email)}.
\end{block}

\begin{block}{Entity Instance}
Một đối tượng cụ thể thuộc entity type, ví dụ: \texttt{S001, Nguyen An, 2005-01-15, an@example.com}.
\end{block}

\begin{block}{Entity Set}
Tập hợp tất cả các entity cùng một entity type, ví dụ tập hợp tất cả sinh viên trong trường.
\end{block}
\end{frame}

%------------------------------------------------
\begin{frame}{Ghi nhớ về Entity Set}
\small
ER Diagram thường biểu diễn \kw{entity type} hoặc \kw{entity set}, không biểu diễn từng dòng dữ liệu cụ thể.

\vspace{2mm}
\begin{exampleblock}{Ví dụ}
ERD biểu diễn cấu trúc \texttt{Student}, chứ không vẽ từng sinh viên cụ thể như S001, S002, S003.
\end{exampleblock}
\end{frame}

%------------------------------------------------
\begin{frame}{Hình minh họa: Entity Set}
\centering
\imageplaceholder{entity-set.png}{Entity type, entity instance và entity set}

\vspace{2mm}
{\small Hình 3. Entity type, entity instance và entity set.}
\end{frame}

%------------------------------------------------
\begin{frame}{Strong Entity}
\small
\kw{Strong Entity} là loại entity có thể được định danh duy nhất bằng chính thuộc tính của nó.

\vspace{2mm}
Đặc điểm:
\begin{itemize}
\smallitem
  \item Có \kw{key attribute} hoặc \kw{primary key}.
  \item Không phụ thuộc vào entity khác để xác định danh tính.
  \item Thường được biểu diễn bằng \kw{hình chữ nhật đơn} trong ER Diagram.
\end{itemize}

\begin{exampleblock}{Ví dụ}
\texttt{Student} có \texttt{student\_id}; \texttt{Employee} có \texttt{employee\_id}; \texttt{Product} có \texttt{product\_id}.
\end{exampleblock}
\end{frame}

%------------------------------------------------
\begin{frame}{Weak Entity}
\small
\kw{Weak Entity} là loại entity không thể được định danh duy nhất chỉ bằng thuộc tính của chính nó.

\vspace{2mm}
Weak Entity cần phụ thuộc vào một \kw{Strong Entity} để được xác định.
\begin{itemize}
\smallitem
  \item Không có khóa riêng đủ mạnh để định danh duy nhất.
  \item Phụ thuộc vào một identifying strong entity.
  \item Thường có \kw{total participation} trong quan hệ với strong entity.
  \item Được biểu diễn bằng \kw{hình chữ nhật kép}.
  \item Quan hệ định danh thường được biểu diễn bằng \kw{hình thoi kép}.
\end{itemize}
\end{frame}

%------------------------------------------------
\begin{frame}{Ví dụ Weak Entity}
\small
Một công ty lưu thông tin người phụ thuộc của nhân viên.

\vspace{2mm}
\begin{itemize}
\smallitem
  \item \texttt{Employee} là strong entity.
  \item \texttt{Dependent} là weak entity.
  \item \texttt{Dependent} không thể tồn tại độc lập nếu không gắn với một nhân viên cụ thể.
\end{itemize}
\end{frame}

%------------------------------------------------
\begin{frame}{Hình minh họa: Strong Entity và Weak Entity}
\centering
\imageplaceholder{strong-weak-entity.png}{Strong Entity và Weak Entity}

\vspace{2mm}
{\small Hình 4. Strong Entity và Weak Entity.}
\end{frame}

%------------------------------------------------
\begin{frame}{Quiz nhanh 2: Entity}
\small
\textbf{Câu 4.} Entity trong ER Model là gì?
\begin{enumerate}
\tightlist
  \qoption{A}{Một kiểu dữ liệu số}
  \qoption{B}{Một loại chỉ mục}
  \qoption{C}{Một câu lệnh SQL}
  \qoption{D}{Một đối tượng hoặc khái niệm trong thế giới thực cần lưu trữ thông tin}
\end{enumerate}

\textbf{Câu 5.} Entity nào sau đây thường là ví dụ của strong entity?
\begin{enumerate}
\tightlist
  \qoption{A}{Một số điện thoại phụ không có chủ sở hữu}
  \qoption{B}{Một dòng log tạm thời không có định danh}
  \qoption{C}{\texttt{Student} có \texttt{student\_id} định danh duy nhất}
  \qoption{D}{Một thuộc tính được suy ra từ ngày sinh}
\end{enumerate}
\end{frame}

%------------------------------------------------
\begin{frame}{Quiz nhanh 2: Entity (tiếp)}
\small
\textbf{Câu 6.} Weak Entity có đặc điểm nào?
\begin{enumerate}
\tightlist
  \qoption{A}{Luôn là một thuộc tính đa trị}
  \qoption{B}{Phụ thuộc vào strong entity để được định danh}
  \qoption{C}{Không thể xuất hiện trong ER Diagram}
  \qoption{D}{Luôn được biểu diễn bằng hình oval đơn}
\end{enumerate}
\end{frame}

%================================================
\section{Attribute trong ER Model}

%------------------------------------------------
\begin{frame}{Attribute là gì?}
\small
\kw{Attribute} là thuộc tính dùng để mô tả entity.

\vspace{2mm}
Ví dụ với entity \texttt{Student}, các attribute có thể gồm:
\begin{itemize}
\smallitem
  \item \texttt{Roll\_No}
  \item \texttt{Name}
  \item \texttt{DOB}
  \item \texttt{Age}
  \item \texttt{Address}
  \item \texttt{Mobile\_No}
\end{itemize}

Trong ER Diagram, attribute thường được biểu diễn bằng \kw{hình oval}.
\end{frame}

%------------------------------------------------
\begin{frame}{Hình minh họa: Attribute}
\centering
\imageplaceholder{attribute.png}{Attribute trong ER Model}
\end{frame}

%------------------------------------------------
\begin{frame}{Các loại Attribute}
\small
Các loại attribute thường gặp:
\begin{enumerate}
\smallitem
  \item \kw{Simple Attribute}
  \item \kw{Key Attribute}
  \item \kw{Composite Attribute}
  \item \kw{Multivalued Attribute}
  \item \kw{Derived Attribute}
\end{enumerate}
\end{frame}

%------------------------------------------------
\begin{frame}{Simple Attribute}
\small
\kw{Simple Attribute} là thuộc tính không thể chia nhỏ hơn thành các thuộc tính có ý nghĩa độc lập.

\vspace{2mm}
Ví dụ:
\begin{itemize}
\smallitem
  \item \texttt{Age}
  \item \texttt{Gender}
  \item \texttt{Roll\_No}
\end{itemize}
\end{frame}

%------------------------------------------------
\begin{frame}{Hình minh họa: Simple Attribute}
\centering
\imageplaceholder{simple-attribute.png}{Simple Attribute trong ER Model}
\end{frame}

%------------------------------------------------
\begin{frame}{Key Attribute}
\small
\kw{Key Attribute} là thuộc tính dùng để định danh duy nhất mỗi entity trong entity set.

\vspace{2mm}
Ví dụ:
\begin{itemize}
\smallitem
  \item \texttt{Roll\_No} định danh duy nhất mỗi sinh viên.
  \item \texttt{Employee\_ID} định danh duy nhất mỗi nhân viên.
  \item \texttt{Product\_ID} định danh duy nhất mỗi sản phẩm.
\end{itemize}

Trong ER Diagram, key attribute được biểu diễn bằng oval có \kw{gạch chân}.
\end{frame}

%------------------------------------------------
\begin{frame}{Hình minh họa: Key Attribute}
\centering
\imageplaceholder{key-attribute.png}{Key Attribute trong ER Model}
\end{frame}

%------------------------------------------------
\begin{frame}{Composite Attribute}
\small
\kw{Composite Attribute} là thuộc tính có thể được chia thành nhiều thuộc tính nhỏ hơn.

\vspace{2mm}
Ví dụ, \texttt{Address} có thể gồm:
\begin{itemize}
\smallitem
  \item \texttt{Street}
  \item \texttt{City}
  \item \texttt{State}
  \item \texttt{Country}
\end{itemize}
\end{frame}

%------------------------------------------------
\begin{frame}{Hình minh họa: Composite Attribute}
\centering
\imageplaceholder{composite-attribute.png}{Composite Attribute trong ER Model}
\end{frame}

%------------------------------------------------
\begin{frame}{Multivalued Attribute}
\small
\kw{Multivalued Attribute} là thuộc tính có thể có nhiều giá trị đối với một entity.

\vspace{2mm}
Ví dụ:
\begin{itemize}
\smallitem
  \item Một sinh viên có thể có nhiều số điện thoại.
  \item Một nhân viên có thể có nhiều kỹ năng.
  \item Một sản phẩm có thể có nhiều màu sắc.
\end{itemize}

Trong ER Diagram, multivalued attribute được biểu diễn bằng \kw{oval kép}.
\end{frame}

%------------------------------------------------
\begin{frame}{Hình minh họa: Multivalued Attribute}
\centering
\imageplaceholder{multivalued-attribute.png}{Multivalued Attribute trong ER Model}
\end{frame}

%------------------------------------------------
\begin{frame}{Derived Attribute}
\small
\kw{Derived Attribute} là thuộc tính có thể được suy ra từ thuộc tính khác.

\vspace{2mm}
Ví dụ:
\begin{itemize}
\smallitem
  \item \texttt{Age} có thể được suy ra từ \texttt{DOB}.
  \item \texttt{Total\_Amount} có thể được tính từ số lượng và đơn giá.
  \item \texttt{Years\_of\_Service} có thể được tính từ ngày bắt đầu làm việc.
\end{itemize}

Trong ER Diagram, derived attribute được biểu diễn bằng \kw{oval nét đứt}.
\end{frame}

%------------------------------------------------
\begin{frame}{Hình minh họa: Derived Attribute}
\centering
\imageplaceholder{derived-attribute.png}{Derived Attribute trong ER Model}
\end{frame}

%------------------------------------------------
\begin{frame}{Quiz nhanh 3: Attribute}
\small
\textbf{Câu 7.} Attribute trong ER Model dùng để làm gì?
\begin{enumerate}
\tightlist
  \qoption{A}{Xóa toàn bộ cơ sở dữ liệu}
  \qoption{B}{Tăng kích thước ổ cứng}
  \qoption{C}{Mô tả thuộc tính của entity}
  \qoption{D}{Thay thế hệ điều hành}
\end{enumerate}

\textbf{Câu 8.} \texttt{Roll\_No} dùng để định danh duy nhất sinh viên là loại attribute nào?
\begin{enumerate}
\tightlist
  \qoption{A}{Composite Attribute}
  \qoption{B}{Multivalued Attribute}
  \qoption{C}{Derived Attribute}
  \qoption{D}{Key Attribute}
\end{enumerate}
\end{frame}

%------------------------------------------------
\begin{frame}{Quiz nhanh 3: Attribute (tiếp)}
\small
\textbf{Câu 9.} \texttt{Address} gồm \texttt{Street}, \texttt{City}, \texttt{State}, \texttt{Country} là ví dụ của loại attribute nào?
\begin{enumerate}
\tightlist
  \qoption{A}{Composite Attribute}
  \qoption{B}{Derived Attribute}
  \qoption{C}{Multivalued Attribute}
  \qoption{D}{Weak Attribute}
\end{enumerate}

\textbf{Câu 10.} \texttt{Age} được tính từ \texttt{DOB} là ví dụ của loại attribute nào?
\begin{enumerate}
\tightlist
  \qoption{A}{Key Attribute}
  \qoption{B}{Derived Attribute}
  \qoption{C}{Composite Attribute}
  \qoption{D}{Multivalued Attribute}
\end{enumerate}
\end{frame}

%================================================
\section{Relationship và Degree}

%------------------------------------------------
\begin{frame}{Relationship trong ER Model}
\small
\kw{Relationship} mô tả mối liên kết giữa các entity.

\vspace{2mm}
Ví dụ:
\begin{itemize}
\smallitem
  \item \texttt{Student} \kw{enrolls in} \texttt{Course}.
  \item \texttt{Employee} \kw{works for} \texttt{Department}.
  \item \texttt{Customer} \kw{places} \texttt{Order}.
  \item \texttt{Doctor} \kw{performs} \texttt{Surgery}.
\end{itemize}

Trong ER Diagram, relationship thường được biểu diễn bằng \kw{hình thoi} và nối với entity bằng các đường liên kết.
\end{frame}

%------------------------------------------------
\begin{frame}{Relationship Type và Relationship Set}
\small
\begin{block}{Relationship Type}
Mô tả loại quan hệ giữa các entity type. Ví dụ: \texttt{Student -- Enrolled in -- Course}.
\end{block}

\begin{block}{Relationship Set}
Tập hợp các relationship cùng loại, ví dụ S1 đăng ký C2, S2 đăng ký C1, S3 đăng ký C3.
\end{block}
\end{frame}

%------------------------------------------------
\begin{frame}{Hình minh họa: Relationship Set}
\centering
\imageplaceholder{relationship-set.png}{Relationship type và relationship set}

\vspace{2mm}
{\small Hình 5. Relationship type và relationship set.}
\end{frame}

%------------------------------------------------
\begin{frame}{Degree of Relationship Set}
\small
\kw{Degree} của relationship set là số lượng entity set khác nhau tham gia vào relationship set.

\vspace{2mm}
Có bốn dạng phổ biến:
\begin{enumerate}
\smallitem
  \item \kw{Unary / Recursive Relationship}: một entity set tham gia.
  \item \kw{Binary Relationship}: hai entity set tham gia.
  \item \kw{Ternary Relationship}: ba entity set tham gia.
  \item \kw{N-ary Relationship}: n entity set tham gia.
\end{enumerate}
\end{frame}

%------------------------------------------------
\begin{frame}{Ví dụ về Degree}
\small
\begin{tabularx}{\textwidth}{>{\bfseries}p{0.28\textwidth} X}
\toprule
Loại degree & Ví dụ \\
\midrule
Unary / Recursive & Một nhân viên quản lý một nhân viên khác \\
Binary & Sinh viên đăng ký môn học \\
Ternary & Nhà cung cấp cung cấp linh kiện cho dự án \\
N-ary & Quan hệ có nhiều hơn ba entity set tham gia \\
\bottomrule
\end{tabularx}
\end{frame}

%------------------------------------------------
\begin{frame}{Hình minh họa: Unary / Recursive Relationship}
\centering
\imageplaceholder{recursive-relationship.png}{Unary hoặc Recursive Relationship}
\end{frame}

%------------------------------------------------
\begin{frame}{Hình minh họa: Binary Relationship}
\centering
\imageplaceholder{binary-relationship.png}{Binary Relationship}
\end{frame}

%------------------------------------------------
\begin{frame}{Hình minh họa: Ternary Relationship}
\centering
\imageplaceholder{ternary-relationship.png}{Ternary Relationship}
\end{frame}

%------------------------------------------------
\begin{frame}{Hình minh họa: N-ary Relationship}
\centering
\imageplaceholder{n-ary-relationship.png}{N-ary Relationship}
\end{frame}

%------------------------------------------------
\begin{frame}{Quiz nhanh 4: Relationship và Degree}
\small
\textbf{Câu 11.} Relationship trong ER Model biểu diễn điều gì?
\begin{enumerate}
\tightlist
  \qoption{A}{Một loại font chữ}
  \qoption{B}{Mối liên kết giữa các entity}
  \qoption{C}{Một chỉ số hiệu năng CPU}
  \qoption{D}{Một kiểu dữ liệu hình ảnh}
\end{enumerate}

\textbf{Câu 12.} Relationship có hai entity set tham gia được gọi là gì?
\begin{enumerate}
\tightlist
  \qoption{A}{Unary Relationship}
  \qoption{B}{Ternary Relationship}
  \qoption{C}{Binary Relationship}
  \qoption{D}{N-ary Relationship}
\end{enumerate}
\end{frame}

%------------------------------------------------
\begin{frame}{Quiz nhanh 4: Relationship và Degree (tiếp)}
\small
\textbf{Câu 13.} Quan hệ một nhân viên quản lý một nhân viên khác là ví dụ của loại degree nào?
\begin{enumerate}
\tightlist
  \qoption{A}{Ternary Relationship}
  \qoption{B}{Binary Relationship}
  \qoption{C}{N-ary Relationship}
  \qoption{D}{Unary / Recursive Relationship}
\end{enumerate}

\textbf{Câu 14.} Degree của relationship set cho biết điều gì?
\begin{enumerate}
\tightlist
  \qoption{A}{Số lượng entity set tham gia vào relationship set}
  \qoption{B}{Số lượng bản ghi trong database}
  \qoption{C}{Số lượng cột trong một bảng}
  \qoption{D}{Số lượng index được tạo ra}
\end{enumerate}
\end{frame}

%================================================
\section{Cardinality trong ER Model}

%------------------------------------------------
\begin{frame}{Cardinality là gì?}
\small
\kw{Cardinality} cho biết số lần tối đa mà một entity trong một entity set có thể tham gia vào một relationship set.

\vspace{2mm}
Cardinality giúp trả lời:
\begin{itemize}
\smallitem
  \item Một entity có thể liên kết với bao nhiêu entity khác?
  \item Quan hệ giữa hai entity là một-một, một-nhiều hay nhiều-nhiều?
  \item Khi chuyển sang bảng quan hệ, cần đặt khóa ngoại hoặc bảng trung gian như thế nào?
\end{itemize}
\end{frame}

%------------------------------------------------
\begin{frame}{One-to-One Relationship}
\small
Quan hệ \kw{one-to-one} xảy ra khi mỗi entity trong mỗi entity set chỉ tham gia tối đa một lần trong relationship.

\vspace{2mm}
Ví dụ:
\begin{itemize}
\smallitem
  \item Một người chỉ có một hộ chiếu.
  \item Một hộ chiếu chỉ thuộc về một người.
\end{itemize}

\begin{center}
\texttt{Person 1 -- 1 Passport}
\end{center}
\end{frame}

%------------------------------------------------
\begin{frame}{Hình minh họa: One-to-One}
\centering
\imageplaceholder{one-to-one.png}{One-to-One Relationship}
\end{frame}

%------------------------------------------------
\begin{frame}{Set representation: One-to-One}
\centering
\imageplaceholder{one-to-one-set-representation.png}{Set Representation of One-to-One Relationship}
\end{frame}

%------------------------------------------------
\begin{frame}{One-to-Many Relationship}
\small
Quan hệ \kw{one-to-many} xảy ra khi một entity ở bên thứ nhất có thể liên kết với nhiều entity ở bên thứ hai.

\vspace{2mm}
Ví dụ:
\begin{itemize}
\smallitem
  \item Một khoa có nhiều giảng viên.
  \item Một khách hàng có nhiều đơn hàng.
  \item Một lớp học có nhiều sinh viên.
\end{itemize}

\begin{center}
\texttt{Department 1 -- M Doctor}
\end{center}
\end{frame}

%------------------------------------------------
\begin{frame}{Hình minh họa: One-to-Many}
\centering
\imageplaceholder{one-to-many.png}{One-to-Many Relationship}
\end{frame}

%------------------------------------------------
\begin{frame}{Set representation: One-to-Many}
\centering
\imageplaceholder{one-to-many-set-representation.png}{Set Representation of One-to-Many Relationship}
\end{frame}

%------------------------------------------------
\begin{frame}{Many-to-One Relationship}
\small
Quan hệ \kw{many-to-one} là chiều ngược lại của one-to-many.

\vspace{2mm}
Ví dụ:
\begin{itemize}
\smallitem
  \item Nhiều ca phẫu thuật có thể do một bác sĩ thực hiện.
  \item Nhiều đơn hàng thuộc về một khách hàng.
  \item Nhiều nhân viên thuộc về một phòng ban.
\end{itemize}

\begin{center}
\texttt{Surgery M -- 1 Surgeon}
\end{center}
\end{frame}

%------------------------------------------------
\begin{frame}{Hình minh họa: Many-to-One}
\centering
\imageplaceholder{many-to-one.png}{Many-to-One Relationship}
\end{frame}

%------------------------------------------------
\begin{frame}{Set representation: Many-to-One}
\centering
\imageplaceholder{many-to-one-set-representation.png}{Set Representation of Many-to-One Relationship}
\end{frame}

%------------------------------------------------
\begin{frame}{Many-to-Many Relationship}
\small
Quan hệ \kw{many-to-many} xảy ra khi entity ở cả hai bên đều có thể tham gia nhiều lần trong relationship.

\vspace{2mm}
Ví dụ:
\begin{itemize}
\smallitem
  \item Một sinh viên học nhiều môn học.
  \item Một môn học có nhiều sinh viên.
  \item Một nhân viên làm nhiều dự án.
  \item Một dự án có nhiều nhân viên.
\end{itemize}
\end{frame}

%------------------------------------------------
\begin{frame}[fragile]{Bảng trung gian cho Many-to-Many}
\small
Khi chuyển sang mô hình quan hệ, many-to-many thường cần một bảng trung gian.

\begin{verbatim}
Students(student_id, full_name)
Courses(course_id, course_name)
Enrollments(student_id, course_id, semester, grade)
\end{verbatim}
\end{frame}

%------------------------------------------------
\begin{frame}{Hình minh họa: Many-to-Many}
\centering
\imageplaceholder{many-to-many.png}{Many-to-Many Relationship}
\end{frame}

%------------------------------------------------
\begin{frame}{Set representation: Many-to-Many}
\centering
\imageplaceholder{many-to-many-set-representation.png}{Set Representation of Many-to-Many Relationship}
\end{frame}

%------------------------------------------------
\begin{frame}{Quiz nhanh 5: Cardinality}
\small
\textbf{Câu 15.} Cardinality trong ER Model cho biết điều gì?
\begin{enumerate}
\tightlist
  \qoption{A}{Số lần tối đa một entity có thể tham gia vào relationship set}
  \qoption{B}{Màu sắc của hình chữ nhật trong ERD}
  \qoption{C}{Số lượng file backup}
  \qoption{D}{Tốc độ truy vấn SQL}
\end{enumerate}

\textbf{Câu 16.} Một người có một hộ chiếu và một hộ chiếu thuộc về một người là quan hệ gì?
\begin{enumerate}
\tightlist
  \qoption{A}{One-to-Many}
  \qoption{B}{Many-to-One}
  \qoption{C}{One-to-One}
  \qoption{D}{Many-to-Many}
\end{enumerate}
\end{frame}

%------------------------------------------------
\begin{frame}{Quiz nhanh 5: Cardinality (tiếp)}
\small
\textbf{Câu 17.} Một khách hàng có nhiều đơn hàng là ví dụ của quan hệ gì?
\begin{enumerate}
\tightlist
  \qoption{A}{One-to-One}
  \qoption{B}{Many-to-Many}
  \qoption{C}{Many-to-One}
  \qoption{D}{One-to-Many}
\end{enumerate}

\textbf{Câu 18.} Quan hệ sinh viên và môn học thường là quan hệ gì nếu một sinh viên học nhiều môn và một môn có nhiều sinh viên?
\begin{enumerate}
\tightlist
  \qoption{A}{One-to-One}
  \qoption{B}{Many-to-Many}
  \qoption{C}{One-to-Many}
  \qoption{D}{Recursive}
\end{enumerate}
\end{frame}

%================================================
\section{Participation Constraint}

%------------------------------------------------
\begin{frame}{Participation Constraint là gì?}
\small
\kw{Participation Constraint} là ràng buộc mô tả việc entity trong một entity set có bắt buộc phải tham gia vào relationship hay không.

\vspace{2mm}
Có hai loại chính:
\begin{enumerate}
\smallitem
  \item \kw{Total Participation}
  \item \kw{Partial Participation}
\end{enumerate}
\end{frame}

%------------------------------------------------
\begin{frame}{Total Participation}
\small
\kw{Total Participation} nghĩa là mọi entity trong entity set đều phải tham gia vào relationship.

\vspace{2mm}
Ví dụ: nếu hệ thống quy định mỗi sinh viên bắt buộc phải đăng ký ít nhất một môn học, thì entity set \texttt{Student} có total participation trong relationship \texttt{Enrolled in}.

\vspace{2mm}
Trong ER Diagram, total participation thường được biểu diễn bằng \kw{đường đôi}.
\end{frame}

%------------------------------------------------
\begin{frame}{Hình minh họa: Total Participation}
\centering
\imageplaceholder{total-participation.png}{Total Participation trong ER Model}
\end{frame}

%------------------------------------------------
\begin{frame}{Partial Participation}
\small
\kw{Partial Participation} nghĩa là entity trong entity set có thể tham gia hoặc không tham gia vào relationship.

\vspace{2mm}
Ví dụ: một số môn học có thể chưa có sinh viên nào đăng ký. Khi đó entity set \texttt{Course} có partial participation trong relationship \texttt{Enrolled in}.

\vspace{2mm}
Trong ER Diagram, partial participation thường được biểu diễn bằng \kw{đường đơn}.
\end{frame}

%------------------------------------------------
\begin{frame}{Set representation: Participation Constraint}
\centering
\imageplaceholder{participation-set-representation.png}{Set Representation of Participation Constraint}
\end{frame}

%------------------------------------------------
\begin{frame}{Quiz nhanh 6: Participation Constraint}
\small
\textbf{Câu 19.} Total Participation có nghĩa là gì?
\begin{enumerate}
\tightlist
  \qoption{A}{Entity không bao giờ tham gia relationship}
  \qoption{B}{Chỉ một entity được tham gia relationship}
  \qoption{C}{Tất cả entity trong entity set bắt buộc tham gia relationship}
  \qoption{D}{Relationship bị xóa khỏi ERD}
\end{enumerate}

\textbf{Câu 20.} Partial Participation có nghĩa là gì?
\begin{enumerate}
\tightlist
  \qoption{A}{Entity có thể tham gia hoặc không tham gia relationship}
  \qoption{B}{Entity bắt buộc phải tham gia relationship}
  \qoption{C}{Entity luôn là weak entity}
  \qoption{D}{Entity luôn có nhiều khóa chính}
\end{enumerate}
\end{frame}

%================================================
\section{Cách vẽ ER Diagram}

%------------------------------------------------
\begin{frame}{Quy trình vẽ ER Diagram}
\small
Có thể vẽ ER Diagram theo sáu bước:
\begin{enumerate}
\smallitem
  \item Xác định entity.
  \item Xác định relationship.
  \item Thêm attribute.
  \item Xác định primary key.
  \item Loại bỏ dư thừa.
  \item Kiểm tra độ rõ ràng.
\end{enumerate}
\end{frame}

%------------------------------------------------
\begin{frame}{Bước 1: Xác định Entity}
\small
Xác định tất cả các entity quan trọng trong hệ thống.

\vspace{2mm}
Ví dụ với hệ thống quản lý sinh viên:
\begin{itemize}
\smallitem
  \item \texttt{Student}
  \item \texttt{Course}
  \item \texttt{Lecturer}
  \item \texttt{Class}
  \item \texttt{Department}
\end{itemize}

Biểu diễn entity bằng \kw{hình chữ nhật}.
\end{frame}

%------------------------------------------------
\begin{frame}{Bước 2: Xác định Relationship}
\small
Xác định quan hệ giữa các entity.

\vspace{2mm}
Ví dụ:
\begin{itemize}
\smallitem
  \item \texttt{Student} đăng ký \texttt{Course}.
  \item \texttt{Lecturer} dạy \texttt{Class}.
  \item \texttt{Department} quản lý \texttt{Lecturer}.
\end{itemize}

Biểu diễn relationship bằng \kw{hình thoi}.
\end{frame}


%------------------------------------------------
\begin{frame}{Bước 3: Thêm Attribute}
\small
Gắn các attribute vào entity bằng các hình oval.

\vspace{2mm}
Ví dụ entity \texttt{Student} có thể có:
\begin{itemize}
\smallitem
  \item \texttt{student\_id}
  \item \texttt{full\_name}
  \item \texttt{date\_of\_birth}
  \item \texttt{email}
  \item \texttt{phone\_number}
\end{itemize}
\end{frame}

%------------------------------------------------
\begin{frame}{Bước 4: Xác định Primary Key}
\small
Với mỗi entity, xác định thuộc tính dùng để định danh duy nhất.

\vspace{2mm}
Ví dụ:
\begin{itemize}
\smallitem
  \item \texttt{Student}: \texttt{student\_id}
  \item \texttt{Course}: \texttt{course\_id}
  \item \texttt{Lecturer}: \texttt{lecturer\_id}
\end{itemize}

Trong ER Diagram, key attribute thường được gạch chân.
\end{frame}
%------------------------------------------------

\begin{frame}{Bước 5 và 6: Hoàn thiện ER Diagram}
\small
\begin{block}{Bước 5: Loại bỏ dư thừa}
Loại bỏ entity không cần thiết, relationship trùng lặp, attribute bị lặp và quan hệ không có ý nghĩa nghiệp vụ.
\end{block}

\begin{block}{Bước 6: Kiểm tra độ rõ ràng}
Kiểm tra sơ đồ có dễ hiểu không, có thể chuyển sang bảng quan hệ không, có phản ánh đúng yêu cầu hệ thống không.
\end{block}
\end{frame}

%================================================
\section{Ví dụ tổng hợp: Hệ thống đăng ký học phần}

%------------------------------------------------
\begin{frame}{Bài toán ví dụ}
\small
Xét hệ thống đăng ký học phần của trường đại học.

\vspace{2mm}
Ta cần xác định:
\begin{itemize}
\smallitem
  \item Entity.
  \item Attribute.
  \item Relationship.
  \item Cardinality.
  \item Cách chuyển sang bảng quan hệ.
\end{itemize}
\end{frame}

%------------------------------------------------
\begin{frame}{Entity chính}
\small
Các entity chính:
\begin{itemize}
\smallitem
  \item \texttt{Student}
  \item \texttt{Course}
  \item \texttt{Lecturer}
  \item \texttt{Department}
\end{itemize}
\end{frame}

%------------------------------------------------
\begin{frame}{Attribute của các entity}
\scriptsize
\begin{columns}[T]
\begin{column}{0.48\textwidth}
\textbf{Student}
\begin{itemize}
\smallitem
  \item \texttt{student\_id}
  \item \texttt{full\_name}
  \item \texttt{email}
  \item \texttt{date\_of\_birth}
\end{itemize}

\textbf{Course}
\begin{itemize}
\smallitem
  \item \texttt{course\_id}
  \item \texttt{course\_name}
  \item \texttt{credits}
\end{itemize}
\end{column}
\begin{column}{0.48\textwidth}
\textbf{Lecturer}
\begin{itemize}
\smallitem
  \item \texttt{lecturer\_id}
  \item \texttt{full\_name}
  \item \texttt{email}
\end{itemize}

\textbf{Department}
\begin{itemize}
\smallitem
  \item \texttt{department\_id}
  \item \texttt{department\_name}
\end{itemize}
\end{column}
\end{columns}
\end{frame}

%------------------------------------------------
\begin{frame}{Relationship chính}
\small
Các relationship chính:
\begin{itemize}
\smallitem
  \item \texttt{Student} \kw{enrolls in} \texttt{Course}.
  \item \texttt{Lecturer} \kw{teaches} \texttt{Course}.
  \item \texttt{Department} \kw{offers} \texttt{Course}.
  \item \texttt{Department} \kw{manages} \texttt{Lecturer}.
\end{itemize}
\end{frame}

%------------------------------------------------
\begin{frame}{Cardinality trong ví dụ}
\small
Một số cardinality:
\begin{itemize}
\smallitem
  \item Một sinh viên có thể đăng ký nhiều học phần.
  \item Một học phần có thể có nhiều sinh viên đăng ký.
  \item Một giảng viên có thể dạy nhiều học phần.
  \item Một học phần có thể do một hoặc nhiều giảng viên phụ trách tùy quy định.
  \item Một khoa có thể quản lý nhiều giảng viên.
\end{itemize}
\end{frame}

%------------------------------------------------
\begin{frame}[fragile]{Chuyển sang bảng quan hệ}
\small
Có thể chuyển thành các bảng:
\begin{verbatim}
Students(student_id, full_name, email, date_of_birth)
Courses(course_id, course_name, credits)
Lecturers(lecturer_id, full_name, email)
Departments(department_id, department_name)
Enrollments(student_id, course_id, semester, grade)
Teaches(lecturer_id, course_id, semester)
\end{verbatim}
\end{frame}

%------------------------------------------------
\begin{frame}{Diễn giải bảng trung gian}
\small
\begin{itemize}
\smallitem
  \item \texttt{student\_id} là khóa chính của \texttt{Students}.
  \item \texttt{course\_id} là khóa chính của \texttt{Courses}.
  \item \texttt{lecturer\_id} là khóa chính của \texttt{Lecturers}.
  \item \texttt{department\_id} là khóa chính của \texttt{Departments}.
  \item \texttt{Enrollments} là bảng trung gian cho quan hệ nhiều-nhiều giữa \texttt{Student} và \texttt{Course}.
  \item \texttt{Teaches} là bảng liên kết giữa \texttt{Lecturer} và \texttt{Course}.
\end{itemize}
\end{frame}

%------------------------------------------------
\begin{frame}{Quiz nhanh 7: Quy trình vẽ ERD}
\small
\textbf{Câu 21.} Khi vẽ ER Diagram, bước đầu tiên thường là gì?
\begin{enumerate}
\tightlist
  \qoption{A}{Chọn màu nền cho sơ đồ}
  \qoption{B}{Tạo index cho bảng}
  \qoption{C}{Xóa toàn bộ relationship}
  \qoption{D}{Xác định các entity quan trọng}
\end{enumerate}

\textbf{Câu 22.} Trong ER Diagram, relationship thường được biểu diễn bằng ký hiệu nào?
\begin{enumerate}
\tightlist
  \qoption{A}{Hình oval}
  \qoption{B}{Hình chữ nhật}
  \qoption{C}{Đường nét đứt}
  \qoption{D}{Hình thoi}
\end{enumerate}
\end{frame}

%================================================
\section{Câu hỏi và bài tập}

%------------------------------------------------
\begin{frame}[allowframebreaks]{Câu hỏi tự luận ngắn}
\small
\begin{enumerate}
\smallitem
  \item ER Model là gì? Vì sao ER Model được xem là mô hình khái niệm trong thiết kế cơ sở dữ liệu?
  \item Phân biệt entity, attribute và relationship.
  \item Phân biệt strong entity và weak entity. Cho ví dụ.
  \item Trình bày bốn loại attribute thường gặp trong ER Model.
  \item Degree của relationship set là gì? Cho ví dụ unary, binary và ternary relationship.
  \item Cardinality là gì? Vì sao cardinality quan trọng khi chuyển ER Model sang mô hình quan hệ?
  \item Phân biệt total participation và partial participation.
\end{enumerate}
\end{frame}

%------------------------------------------------
\begin{frame}{Bài tập 1: Hệ thống trường đại học}
\small
Một trường đại học cần xây dựng hệ thống quản lý sinh viên, học phần, giảng viên và khoa.

\vspace{2mm}
\begin{enumerate}
\smallitem
  \item Xác định các entity chính.
  \item Liệt kê một số attribute cho từng entity.
  \item Xác định relationship giữa các entity.
  \item Xác định cardinality cho từng relationship.
  \item Chuyển mô hình ER sang danh sách bảng quan hệ.
\end{enumerate}
\end{frame}

%------------------------------------------------
\begin{frame}{Bài tập 2: Cửa hàng online}
\small
Một cửa hàng online cần quản lý khách hàng, sản phẩm, đơn hàng và thanh toán.

\vspace{2mm}
\begin{enumerate}
\smallitem
  \item Xác định entity và attribute chính.
  \item Xác định relationship giữa \texttt{Customer}, \texttt{Order}, \texttt{Product} và \texttt{Payment}.
  \item Xác định relationship nào là one-to-many, relationship nào là many-to-many.
  \item Đề xuất bảng trung gian nếu cần.
\end{enumerate}
\end{frame}

%------------------------------------------------
\begin{frame}{Bài tập 3: Nhân viên và người phụ thuộc}
\small
Một công ty muốn lưu thông tin nhân viên và người phụ thuộc của nhân viên.

\vspace{2mm}
\begin{enumerate}
\smallitem
  \item Xác định strong entity và weak entity.
  \item Giải thích vì sao một entity là weak entity.
  \item Xác định identifying relationship.
  \item Gợi ý cách chuyển sang bảng quan hệ.
\end{enumerate}
\end{frame}

%------------------------------------------------
\begin{frame}{Bài tập 4: Hệ thống bệnh viện}
\small
Một bệnh viện cần quản lý bệnh nhân, bác sĩ, khoa, ca phẫu thuật và thuốc.

\vspace{2mm}
\begin{enumerate}
\smallitem
  \item Xác định ít nhất 5 entity.
  \item Xác định ít nhất 5 relationship.
  \item Phân loại một số relationship theo degree và cardinality.
  \item Chỉ ra relationship nào có thể cần bảng trung gian.
\end{enumerate}
\end{frame}

%================================================
\section{Tóm tắt và từ khóa}

%------------------------------------------------
\begin{frame}{Tóm tắt bài học}
\small
\begin{itemize}
\smallitem
  \item ER Model là mô hình khái niệm dùng để thiết kế cơ sở dữ liệu.
  \item ER Diagram là biểu diễn đồ họa của ER Model.
  \item Ba thành phần chính của ER Model là entity, attribute và relationship.
  \item Entity là đối tượng hoặc khái niệm cần lưu trữ dữ liệu.
  \item Strong entity có khóa định danh riêng; weak entity phụ thuộc vào strong entity.
  \item Attribute mô tả entity và có nhiều loại như simple, key, composite, multivalued và derived attribute.
\end{itemize}
\end{frame}

%------------------------------------------------
\begin{frame}{Tóm tắt bài học (tiếp)}
\small
\begin{itemize}
\smallitem
  \item Relationship biểu diễn mối liên kết giữa các entity.
  \item Degree cho biết số entity set tham gia vào relationship.
  \item Cardinality cho biết số lần tối đa entity có thể tham gia vào relationship.
  \item Participation constraint cho biết entity tham gia bắt buộc hay tùy chọn vào relationship.
  \item Quy trình vẽ ER Diagram gồm xác định entity, relationship, attribute, khóa chính, loại bỏ dư thừa và kiểm tra độ rõ ràng.
\end{itemize}
\end{frame}

%------------------------------------------------
\begin{frame}{Từ khóa chính}
\scriptsize
\begin{columns}[T]
\begin{column}{0.33\textwidth}
\begin{itemize}
\smallitem
  \item Entity-Relationship Model
  \item ER Model
  \item Entity-Relationship Diagram
  \item ERD
  \item Entity
  \item Entity Type
  \item Entity Set
  \item Entity Instance
  \item Strong Entity
  \item Weak Entity
\end{itemize}
\end{column}
\begin{column}{0.33\textwidth}
\begin{itemize}
\smallitem
  \item Attribute
  \item Simple Attribute
  \item Key Attribute
  \item Composite Attribute
  \item Multivalued Attribute
  \item Derived Attribute
  \item Relationship
  \item Relationship Type
  \item Relationship Set
  \item Degree
\end{itemize}
\end{column}
\begin{column}{0.33\textwidth}
\begin{itemize}
\smallitem
  \item Unary Relationship
  \item Binary Relationship
  \item Ternary Relationship
  \item N-ary Relationship
  \item Cardinality
  \item One-to-One
  \item One-to-Many
  \item Many-to-Many
  \item Participation Constraint
  \item Total/Partial Participation
\end{itemize}
\end{column}
\end{columns}
\end{frame}

%------------------------------------------------
\begin{frame}
\centering
\Huge Cảm ơn!\\[0.5cm]
\Large Câu hỏi và thảo luận
\end{frame}

\end{document}
